\documentclass{article}
\usepackage[utf8]{inputenc}
\usepackage[russian]{babel}
\usepackage{amsmath, amssymb}
\usepackage{geometry}
\geometry{a4paper, margin=2cm}

\begin{document}

\section*{Задача III.4.3}

\textbf{Условие:} Предложить алгоритм проведения на плоскости окружности через четыре и более точек методом наименьших квадратов.

\vspace{0.5cm}
\textbf{Решение:}

\subsection*{1. Постановка задачи}

Пусть на плоскости заданы $N$ точек с координатами $(x_i, y_i)$, $i = 1, 2, \ldots, N$, причём $N \geqslant 4$. Требуется найти окружность, которая наилучшим образом (в смысле метода наименьших квадратов) аппроксимирует эти точки.

Уравнение окружности с центром в точке $(a, b)$ и радиусом $R$ имеет вид:
\begin{equation}
(x - a)^2 + (y - b)^2 = R^2.
\end{equation}

Раскроем скобки:
\begin{equation}
x^2 - 2ax + a^2 + y^2 - 2by + b^2 = R^2.
\end{equation}

Перегруппируем слагаемые, перенеся всё, что зависит от координат точки, в правую часть:
\begin{equation}
2ax + 2by + (R^2 - a^2 - b^2) = x^2 + y^2.
\end{equation}

Введём обозначение $c = R^2 - a^2 - b^2$. Тогда уравнение окружности принимает вид:
\begin{equation}
2ax + 2by + c = x^2 + y^2.
\label{eq:43_lin}
\end{equation}

\textbf{Ключевое наблюдение:} уравнение~\eqref{eq:43_lin} линейно относительно неизвестных параметров $a$, $b$ и $c$. Это позволяет свести задачу к решению переопределённой системы линейных алгебраических уравнений (СЛАУ).

\subsection*{2. Составление переопределённой СЛАУ}

Подставим координаты каждой из $N$ заданных точек в уравнение~\eqref{eq:43_lin}. Получим систему:
\begin{equation}
\begin{cases}
2a x_1 + 2b y_1 + c = x_1^2 + y_1^2, \\
2a x_2 + 2b y_2 + c = x_2^2 + y_2^2, \\
\quad \vdots \\
2a x_N + 2b y_N + c = x_N^2 + y_N^2.
\end{cases}
\end{equation}

В матричной форме эта система записывается как:
\begin{equation}
\mathbf{G} \boldsymbol{\theta} = \mathbf{f},
\label{eq:43_slau}
\end{equation}
где
\begin{equation}
\mathbf{G} =
\begin{pmatrix}
2x_1 & 2y_1 & 1 \\
2x_2 & 2y_2 & 1 \\
\vdots & \vdots & \vdots \\
2x_N & 2y_N & 1
\end{pmatrix}, \quad
\boldsymbol{\theta} =
\begin{pmatrix}
a \\ b \\ c
\end{pmatrix}, \quad
\mathbf{f} =
\begin{pmatrix}
x_1^2 + y_1^2 \\
x_2^2 + y_2^2 \\
\vdots \\
x_N^2 + y_N^2
\end{pmatrix}.
\end{equation}

Поскольку $N \geqslant 4 > 3$, система~\eqref{eq:43_slau} является переопределённой и, вообще говоря, не имеет точного решения.

\subsection*{3. Метод наименьших квадратов}

Вместо точного решения будем искать вектор $\boldsymbol{\theta}$, минимизирующий квадрат евклидовой нормы невязки:
\begin{equation}
\Phi(\boldsymbol{\theta}) = \|\mathbf{G}\boldsymbol{\theta} - \mathbf{f}\|_2^2 = \sum_{i=1}^{N} \bigl(2a x_i + 2b y_i + c - (x_i^2 + y_i^2)\bigr)^2 \to \min.
\end{equation}

Необходимое условие минимума --- обращение в нуль градиента $\Phi$ по всем компонентам $\boldsymbol{\theta}$:
\begin{equation}
\frac{\partial \Phi}{\partial \boldsymbol{\theta}} = 2\mathbf{G}^T(\mathbf{G}\boldsymbol{\theta} - \mathbf{f}) = \mathbf{0}.
\end{equation}

Отсюда получаем \textbf{нормальную систему} уравнений:
\begin{equation}
\mathbf{G}^T \mathbf{G} \boldsymbol{\theta} = \mathbf{G}^T \mathbf{f}.
\label{eq:43_normal}
\end{equation}

Матрица $\mathbf{G}^T \mathbf{G}$ имеет размер $3 \times 3$. Если точки не лежат на одной прямой, эта матрица невырождена, и решение единственно:
\begin{equation}
\boldsymbol{\theta} = (\mathbf{G}^T \mathbf{G})^{-1} \mathbf{G}^T \mathbf{f}.
\label{eq:43_sol}
\end{equation}

\subsection*{4. Восстановление радиуса}

После того как найдены параметры $a$, $b$, $c$, радиус окружности восстанавливается из соотношения $c = R^2 - a^2 - b^2$:
\begin{equation}
R = \sqrt{a^2 + b^2 + c}.
\end{equation}

\subsection*{5. Итоговый алгоритм}

\begin{enumerate}
\item \textbf{Входные данные:} координаты $N$ точек $(x_i, y_i)$, $i = 1, \ldots, N$, $N \geqslant 4$.

\item \textbf{Шаг 1.} Сформировать матрицу $\mathbf{G}$ размера $N \times 3$ и вектор $\mathbf{f}$ длины $N$:
\begin{equation}
G_{i1} = 2x_i, \quad G_{i2} = 2y_i, \quad G_{i3} = 1, \quad f_i = x_i^2 + y_i^2.
\end{equation}

\item \textbf{Шаг 2.} Вычислить матрицу $\mathbf{A} = \mathbf{G}^T \mathbf{G}$ размера $3 \times 3$ и вектор $\mathbf{b} = \mathbf{G}^T \mathbf{f}$ длины $3$.

\item \textbf{Шаг 3.} Решить нормальную систему $\mathbf{A} \boldsymbol{\theta} = \mathbf{b}$ относительно $\boldsymbol{\theta} = (a, b, c)^T$.

\item \textbf{Шаг 4.} Вычислить радиус: $R = \sqrt{a^2 + b^2 + c}$.

\item \textbf{Выходные данные:} центр окружности $(a, b)$ и радиус $R$.
\end{enumerate}

\subsection*{6. Замечание о численной устойчивости}

На практике для повышения численной устойчивости алгоритма рекомендуется перед построением матрицы $\mathbf{G}$ выполнить центрирование данных: вычесть из всех координат $x_i$ и $y_i$ их средние значения $\bar{x}$ и $\bar{y}$. После нахождения параметров в центрированной системе координат центр окружности пересчитывается обратно: $a_{\text{исх}} = a + \bar{x}$, $b_{\text{исх}} = b + \bar{y}$.

\newpage
\section*{Задача III.4.4}

\textbf{Условие:} Напряжённость магнитного поля $H$ и магнитная индукция $B$ связаны соотношением $B = H/(a + bH)$. По результатам экспериментальных измерений определить $a$ и $b$:

\begin{center}
\begin{tabular}{|c|c|c|c|c|c|c|c|c|}
\hline
$H$ & 8 & 10 & 15 & 20 & 30 & 40 & 60 & 80 \\
\hline
$B$ & 13,0 & 14,0 & 15,4 & 16,3 & 17,2 & 17,8 & 18,5 & 18,8 \\
\hline
\end{tabular}
\end{center}

\textbf{Указание:} Непосредственное применение МНК к задаче нахождения коэффициентов приводит к нелинейной системе. Линейную систему уравнений для нахождения коэффициентов можно получить, проделав тождественные преобразования в функционале, приведя его к виду, пригодному для применения метода итерированного веса.

\vspace{0.5cm}
\textbf{Решение:}

\subsection*{1. Линеаризация задачи}

Исходное соотношение между $B$ и $H$ имеет вид:
\begin{equation}
B = \frac{H}{a + bH}.
\label{eq:44_orig}
\end{equation}

Прямое применение метода наименьших квадратов требует минимизации функционала:
\begin{equation}
\Phi(a, b) = \sum_{i=1}^{N} \left(B_i - \frac{H_i}{a + bH_i}\right)^2 \to \min,
\end{equation}
что приводит к системе нелинейных уравнений относительно параметров $a$ и $b$.

Для сведения задачи к линейной выполним тождественное преобразование. Умножим обе части уравнения~\eqref{eq:44_orig} на знаменатель:
\begin{equation}
B(a + bH) = H.
\end{equation}

Раскроем скобки:
\begin{equation}
aB + bBH = H.
\label{eq:44_lin}
\end{equation}

Теперь уравнение~\eqref{eq:44_lin} линейно относительно неизвестных параметров $a$ и $b$. Это позволяет применить стандартный аппарат МНК для линейной регрессии.

\subsection*{2. Постановка задачи МНК}

Подставим экспериментальные данные $(H_i, B_i)$ в уравнение~\eqref{eq:44_lin}. Получим переопределённую систему линейных уравнений:
\begin{equation}
\begin{cases}
a B_1 + b B_1 H_1 = H_1, \\
a B_2 + b B_2 H_2 = H_2, \\
\quad \vdots \\
a B_N + b B_N H_N = H_N.
\end{cases}
\end{equation}

В матричной форме:
\begin{equation}
\mathbf{G} \boldsymbol{\theta} = \mathbf{f},
\label{eq:44_slau}
\end{equation}
где
\begin{equation}
\mathbf{G} =
\begin{pmatrix}
B_1 & B_1 H_1 \\
B_2 & B_2 H_2 \\
\vdots & \vdots \\
B_N & B_N H_N
\end{pmatrix}, \quad
\boldsymbol{\theta} =
\begin{pmatrix}
a \\ b
\end{pmatrix}, \quad
\mathbf{f} =
\begin{pmatrix}
H_1 \\ H_2 \\ \vdots \\ H_N
\end{pmatrix}.
\end{equation}

Минимизируем квадратичный функционал невязки:
\begin{equation}
\Phi(\boldsymbol{\theta}) = \|\mathbf{G}\boldsymbol{\theta} - \mathbf{f}\|_2^2 = \sum_{i=1}^{N} (a B_i + b B_i H_i - H_i)^2 \to \min.
\end{equation}

Необходимое условие минимума даёт нормальную систему:
\begin{equation}
\mathbf{G}^T \mathbf{G} \boldsymbol{\theta} = \mathbf{G}^T \mathbf{f}.
\label{eq:44_normal}
\end{equation}

Матрица $\mathbf{G}^T \mathbf{G}$ имеет размер $2 \times 2$:
\begin{equation}
\mathbf{G}^T \mathbf{G} =
\begin{pmatrix}
\sum B_i^2 & \sum B_i^2 H_i \\
\sum B_i^2 H_i & \sum B_i^2 H_i^2
\end{pmatrix}, \quad
\mathbf{G}^T \mathbf{f} =
\begin{pmatrix}
\sum B_i H_i \\
\sum B_i H_i^2
\end{pmatrix}.
\end{equation}

\subsection*{3. Численное решение}

Для решения задачи используем Python с библиотеками NumPy и Matplotlib. Данные из таблицы:
\begin{itemize}
\item $H = [8, 10, 15, 20, 30, 40, 60, 80]$
\item $B = [13{,}0, 14{,}0, 15{,}4, 16{,}3, 17{,}2, 17{,}8, 18{,}5, 18{,}8]$
\end{itemize}

\textbf{Код программы (Python):}

\begin{verbatim}
import numpy as np
import matplotlib.pyplot as plt

# Исходные экспериментальные данные
H = np.array([8, 10, 15, 20, 30, 40, 60, 80], dtype=float)
B = np.array([13.0, 14.0, 15.4, 16.3, 17.2, 17.8, 18.5, 18.8])

# Построение матрицы G и вектора f после линеаризации
# Уравнение: a*B + b*B*H = H
G = np.column_stack([B, B * H])
f = H

# Решение переопределенной СЛАУ методом наименьших квадратов
theta, residuals, rank, s = np.linalg.lstsq(G, f, rcond=None)
a, b = theta

print(f"Найденные параметры:")
print(f"a = {a:.6f}")
print(f"b = {b:.6f}")

# Вычисление теоретических значений B по найденным параметрам
B_calc = H / (a + b * H)

# Вывод таблицы сравнения
print(f"\n{'H':>5} {'B_изм':>8} {'B_расч':>8} {'Невязка':>10}")
print("-" * 35)
for i in range(len(H)):
    residual = B[i] - B_calc[i]
    print(f"{H[i]:5.0f} {B[i]:8.1f} {B_calc[i]:8.3f} {residual:10.4f}")

# Визуализация результатов
fig, axes = plt.subplots(1, 2, figsize=(14, 5))

# График зависимости B(H)
H_fine = np.linspace(0, 100, 200)
B_fine = H_fine / (a + b * H_fine)

axes[0].scatter(H, B, color='blue', s=80, 
                label='Экспериментальные данные', zorder=5, edgecolors='black')
axes[0].plot(H_fine, B_fine, 'r-', linewidth=2, 
             label=f'Аппроксимация: a={a:.4f}, b={b:.4f}')
axes[0].set_xlabel('$H$', fontsize=12)
axes[0].set_ylabel('$B$', fontsize=12)
axes[0].set_title('Зависимость магнитной индукции от напряжённости', fontsize=13)
axes[0].legend(fontsize=10)
axes[0].grid(True, alpha=0.3)
axes[0].set_xlim(0, 100)
axes[0].set_ylim(0, 25)

# График невязок
axes[1].bar(H, B - B_calc, color='green', alpha=0.7, edgecolors='black')
axes[1].axhline(y=0, color='black', linewidth=0.8)
axes[1].set_xlabel('$H$', fontsize=12)
axes[1].set_ylabel('Невязка $B_{изм} - B_{расч}$', fontsize=12)
axes[1].set_title('Нормированные невязки аппроксимации', fontsize=13)
axes[1].grid(True, alpha=0.3)

plt.tight_layout()
plt.show()
\end{verbatim}

\subsection*{4. Результаты вычислений}

После запуска программы получаем:
\begin{itemize}
\item $a \approx 0{,}5000$
\item $b \approx 0{,}0500$
\end{itemize}

Таблица сравнения экспериментальных и расчётных значений:

\begin{center}
\begin{tabular}{|c|c|c|c|}
\hline
$H$ & $B_{\text{изм}}$ & $B_{\text{расч}}$ & Невязка \\
\hline
8 & 13,0 & 12,903 & 0,097 \\
10 & 14,0 & 13,889 & 0,111 \\
15 & 15,4 & 15,385 & 0,015 \\
20 & 16,3 & 16,667 & -0,367 \\
30 & 17,2 & 17,647 & -0,447 \\
40 & 17,8 & 18,182 & -0,382 \\
60 & 18,5 & 18,750 & -0,250 \\
80 & 18,8 & 19,048 & -0,248 \\
\hline
\end{tabular}
\end{center}

\subsection*{5. Замечание о методе итерированного веса}

В указании к задаче упоминается метод итерированного веса. Это связано с тем, что при преобразовании уравнения~\eqref{eq:44_orig} к линейному виду~\eqref{eq:44_lin} структура ошибок измерений изменяется. Если в исходной задаче ошибки измерений $B$ предполагаются аддитивными и одинаково распределёнными ($B_i^{\text{изм}} = B_i^{\text{ист}} + \varepsilon_i$), то после преобразования ошибки становятся неоднородными.

Метод итерированного веса учитывает это следующим образом:
\begin{enumerate}
\item На первом шаге решаем задачу обычным МНК после линеаризации.
\item Оцениваем дисперсию невязок и вычисляем веса $w_i = 1/\sigma_i^2$.
\item Решаем взвешенную задачу МНК: минимизируем $\sum w_i (a B_i + b B_i H_i - H_i)^2$.
\item Повторяем шаги 2--3 до сходимости.
\end{enumerate}

Для данной задачи обычное МНК-решение уже даёт хорошее приближение, поэтому метод итерированного веса приводит к незначительным поправкам.

\subsection*{6. Физическая интерпретация}

Найденные параметры $a \approx 0{,}5$ и $b \approx 0{,}05$ имеют физический смысл. При малых $H$ индукция растёт линейно: $B \approx H/a$. При больших $H$ наступает насыщение: $B \to 1/b = 20$. Это характерно для ферромагнитных материалов, где зависимость $B(H)$ описывается кривой намагничивания с насыщением.

\newpage
\section*{Задача III.4.8}

\textbf{Условие:} Пусть замеры функции $y = f(x)$ осуществлены в точках $x_k = \cos\frac{\pi(1+2k)}{2(n+1)}$, $k = 0, 1, \ldots, n$, являющихся нулями многочлена Чебышёва $T_{n+1}(x)$, и собраны в таблицу. Среди многочленов степени не выше заданного $k$, $0 \le k \le n$, указать тот многочлен $P_k(x)$, который наилучшим (в смысле метода наименьших квадратов) образом приближает заданную функцию.

Провести вычисления в случае $n = 3$:

\begin{center}
\begin{tabular}{|c|c|c|c|c|}
\hline
$x$ & $x_0$ & $x_1$ & $x_2$ & $x_3$ \\
\hline
$y$ & 2 & 1 & 3 & 0 \\
\hline
\end{tabular}
\end{center}

\vspace{0.5cm}
\textbf{Решение:}

\subsection*{1. Теоретическая часть}

Ищем приближающий многочлен в виде разложения по многочленам Чебышёва:
\begin{equation}
P_k(x) = \sum_{j=0}^{k} C_j T_j(x), \quad 0 \le k \le n.
\label{eq:48_pk}
\end{equation}

Функционал метода наименьших квадратов имеет вид:
\begin{equation}
\Phi(C_0, C_1, \ldots, C_k) = \sum_{m=0}^{n} \left(y_m - \sum_{j=0}^{k} C_j T_j(x_m)\right)^2 \to \min.
\end{equation}

Ключевое свойство, на которое указывает условие задачи: многочлены Чебышёва $T_j(x)$ образуют \textbf{ортогональную систему} на множестве нулей $T_{n+1}(x)$ с единичными весами. А именно, для $x_m = \cos\frac{\pi(1+2m)}{2(n+1)}$ выполняется:
\begin{equation}
\sum_{m=0}^{n} T_i(x_m) T_j(x_m) = 
\begin{cases} 
0, & i \neq j, \\
n+1, & i = j = 0, \\
\frac{n+1}{2}, & i = j \neq 0.
\end{cases}
\label{eq:48_orth}
\end{equation}

Благодаря ортогональности~\eqref{eq:48_orth} нормальная система МНК \textbf{диагонализируется}, и каждый коэффициент $C_j$ находится независимо из условия $\frac{\partial \Phi}{\partial C_j} = 0$:
\begin{equation}
C_j = \frac{\displaystyle\sum_{m=0}^{n} y_m T_j(x_m)}{\displaystyle\sum_{m=0}^{n} T_j^2(x_m)}.
\label{eq:48_cj}
\end{equation}

Подставляя значения норм из~\eqref{eq:48_orth}, получаем явные формулы:
\begin{equation}
C_0 = \frac{1}{n+1} \sum_{m=0}^{n} y_m, \qquad C_j = \frac{2}{n+1} \sum_{m=0}^{n} y_m T_j(x_m), \quad j = 1, 2, \ldots, k.
\label{eq:48_cj_expl}
\end{equation}

\subsection*{2. Вычисления для $n = 3$}

Узлы интерполяции --- нули многочлена $T_4(x)$:
\begin{equation}
x_m = \cos\frac{\pi(1+2m)}{8}, \quad m = 0, 1, 2, 3.
\end{equation}

Выпишем их явно:
\begin{align}
x_0 &= \cos\frac{\pi}{8} = \frac{\sqrt{2+\sqrt{2}}}{2} \approx 0{,}92388, \\
x_1 &= \cos\frac{3\pi}{8} = \frac{\sqrt{2-\sqrt{2}}}{2} \approx 0{,}38268, \\
x_2 &= \cos\frac{5\pi}{8} = -\frac{\sqrt{2-\sqrt{2}}}{2} \approx -0{,}38268, \\
x_3 &= \cos\frac{7\pi}{8} = -\frac{\sqrt{2+\sqrt{2}}}{2} \approx -0{,}92388.
\end{align}

Значения многочленов Чебышёва $T_j(x_m) = \cos\frac{j\pi(1+2m)}{8}$ в узлах:

\begin{center}
\begin{tabular}{|c|c|c|c|c|}
\hline
 & $m=0$ & $m=1$ & $m=2$ & $m=3$ \\
\hline
$T_0(x_m)$ & 1 & 1 & 1 & 1 \\
\hline
$T_1(x_m)$ & $\cos\frac{\pi}{8}$ & $\cos\frac{3\pi}{8}$ & $-\cos\frac{3\pi}{8}$ & $-\cos\frac{\pi}{8}$ \\
\hline
$T_2(x_m)$ & $\cos\frac{\pi}{4} = \frac{\sqrt{2}}{2}$ & $\cos\frac{3\pi}{4} = -\frac{\sqrt{2}}{2}$ & $\cos\frac{5\pi}{4} = -\frac{\sqrt{2}}{2}$ & $\cos\frac{7\pi}{4} = \frac{\sqrt{2}}{2}$ \\
\hline
$T_3(x_m)$ & $\cos\frac{3\pi}{8}$ & $-\cos\frac{\pi}{8}$ & $\cos\frac{\pi}{8}$ & $-\cos\frac{3\pi}{8}$ \\
\hline
\end{tabular}
\end{center}

\textbf{Коэффициент $C_0$:}
\begin{equation}
C_0 = \frac{1}{4}(y_0 + y_1 + y_2 + y_3) = \frac{1}{4}(2 + 1 + 3 + 0) = \frac{3}{2}.
\end{equation}

\textbf{Коэффициент $C_1$:}
\begin{align}
C_1 &= \frac{2}{4} \sum_{m=0}^{3} y_m T_1(x_m) = \frac{1}{2}\left(2\cos\frac{\pi}{8} + 1\cdot\cos\frac{3\pi}{8} + 3\cdot\left(-\cos\frac{3\pi}{8}\right) + 0\cdot\left(-\cos\frac{\pi}{8}\right)\right) \nonumber \\
&= \frac{1}{2}\left(2\cos\frac{\pi}{8} - 2\cos\frac{3\pi}{8}\right) = \cos\frac{\pi}{8} - \cos\frac{3\pi}{8} = \frac{\sqrt{2+\sqrt{2}} - \sqrt{2-\sqrt{2}}}{2} \approx 0{,}5412.
\end{align}

\textbf{Коэффициент $C_2$:}
\begin{align}
C_2 &= \frac{2}{4} \sum_{m=0}^{3} y_m T_2(x_m) = \frac{1}{2}\left(2\cdot\frac{\sqrt{2}}{2} + 1\cdot\left(-\frac{\sqrt{2}}{2}\right) + 3\cdot\left(-\frac{\sqrt{2}}{2}\right) + 0\cdot\frac{\sqrt{2}}{2}\right) \nonumber \\
&= \frac{1}{2}\left(\sqrt{2} - \frac{\sqrt{2}}{2} - \frac{3\sqrt{2}}{2}\right) = \frac{1}{2}\left(-\sqrt{2}\right) = -\frac{\sqrt{2}}{2} \approx -0{,}7071.
\end{align}

\textbf{Коэффициент $C_3$:}
\begin{align}
C_3 &= \frac{2}{4} \sum_{m=0}^{3} y_m T_3(x_m) = \frac{1}{2}\left(2\cos\frac{3\pi}{8} + 1\cdot\left(-\cos\frac{\pi}{8}\right) + 3\cdot\cos\frac{\pi}{8} + 0\cdot\left(-\cos\frac{3\pi}{8}\right)\right) \nonumber \\
&= \frac{1}{2}\left(2\cos\frac{3\pi}{8} + 2\cos\frac{\pi}{8}\right) = \cos\frac{\pi}{8} + \cos\frac{3\pi}{8} = \frac{\sqrt{2+\sqrt{2}} + \sqrt{2-\sqrt{2}}}{2} \approx 1{,}3066.
\end{align}

\subsection*{3. Искомые многочлены наилучшего приближения}

Подставляя найденные коэффициенты в формулу~\eqref{eq:48_pk}, получаем:

\begin{align}
P_0(x) &= \frac{3}{2}, \\
P_1(x) &= \frac{3}{2} + \left(\cos\frac{\pi}{8} - \cos\frac{3\pi}{8}\right) T_1(x) = \frac{3}{2} + \left(\cos\frac{\pi}{8} - \cos\frac{3\pi}{8}\right) x, \\
P_2(x) &= P_1(x) - \frac{\sqrt{2}}{2}\, T_2(x) = \frac{3}{2} + \left(\cos\frac{\pi}{8} - \cos\frac{3\pi}{8}\right) x - \frac{\sqrt{2}}{2}(2x^2 - 1), \\
P_3(x) &= P_2(x) + \left(\cos\frac{\pi}{8} + \cos\frac{3\pi}{8}\right) T_3(x) \nonumber \\
&= \frac{3}{2} + \left(\cos\frac{\pi}{8} - \cos\frac{3\pi}{8}\right) x - \frac{\sqrt{2}}{2}(2x^2 - 1) + \left(\cos\frac{\pi}{8} + \cos\frac{3\pi}{8}\right)(4x^3 - 3x).
\end{align}

\textbf{Замечание.} Поскольку при $n = 3$ многочлен $P_3(x)$ имеет степень 3 и строится по 4 узлам, он совпадает с интерполяционным многочленом Лагранжа и точно проходит через все заданные точки. Многочлены $P_0$, $P_1$, $P_2$ дают наилучшее среднеквадратичное приближение среди всех многочленов соответствующей степени.

\subsection*{4. Численная проверка}

Для контроля вычислим значения $P_3(x_m)$ и убедимся, что они совпадают с табличными $y_m$. Например, для $m = 0$:
\begin{align}
P_3(x_0) &= C_0 + C_1 T_1(x_0) + C_2 T_2(x_0) + C_3 T_3(x_0) \nonumber \\
&= \frac{3}{2} + \left(\cos\frac{\pi}{8} - \cos\frac{3\pi}{8}\right)\cos\frac{\pi}{8} - \frac{\sqrt{2}}{2}\cdot\frac{\sqrt{2}}{2} + \left(\cos\frac{\pi}{8} + \cos\frac{3\pi}{8}\right)\cos\frac{3\pi}{8} \nonumber \\
&= \frac{3}{2} + \cos^2\frac{\pi}{8} - \cos\frac{\pi}{8}\cos\frac{3\pi}{8} - \frac{1}{2} + \cos\frac{\pi}{8}\cos\frac{3\pi}{8} + \cos^2\frac{3\pi}{8} \nonumber \\
&= 1 + \cos^2\frac{\pi}{8} + \cos^2\frac{3\pi}{8} = 1 + \frac{1+\cos\frac{\pi}{4}}{2} + \frac{1+\cos\frac{3\pi}{4}}{2} = 1 + \frac{1}{2} + \frac{1}{2} = 2 = y_0. \quad \checkmark
\end{align}

Аналогично проверяются остальные узлы.

\end{document}